\section{Dataset}
Nell'ambito del machine learning e del deep learning per l'apprendimento di un modello è necessario fare uso di dati di training e di test che saranno prelevati da vari dataset reperibili online. Per questo progetto sono stati utilizzati i seguenti dataset:

\subsection{PhonItaliaR}
PhonItaliaR fornisce rappresentazioni fonologiche per 120.000 forme di parole italiane. Ognuno di questi include anche confini sillabici, marcature di accento e una gamma completa di statistiche lessicali.
Utilizzando dati derivati da questo lessico, è anche possibile generare un insieme di database derivati e fornito stime dell'uso della frequenza posizionale per fonemi, sillabe, inizi e coda sillabiche italiane, nonché bigrammi di caratteri e fonemi.

\Fonte{springer-2013}


\subsection{Q2Stress}
Q2Stress è un database per l'italiano che contiene informazioni su molteplici segnali che i lettori possono
utilizzare per assegnare l'accento. In particolare, il database offre misure di frequenza di tipo e token dei pattern di accento in funzione del numero 
di sillabe, della categoria grammaticale, dell'inizio delle parole, delle terminazioni delle parole e delle strutture consonanti-vocali. Sono disponibili dati 
sia per adulti, sia per i bambini.
Q2Stress può aiutare i ricercatori a rispondere a domande empiriche e teoriche sull'assegnazione dell'accento e sulla relazione tra ortografia e fonologia 
nella lettura. È progettato come una risorsa facile da usare, poiché include script che permettono ai ricercatori di esplorare e selezionare i propri 
stimoli secondo diversi criteri, oltre a tabelle riassuntive per l'analisi complessiva dei dati.

\Fonte{istc-q2stress}


 \subsection{NRC Emotion Lexicon}
 Il NRC Emotion Lexicon (o Word-Emotion Association Lexicon) è un ampio database lessicale per la Sentiment Analysis sviluppato dal National Research Council 
 Canada. È un dataset che associa un vasto vocabolario (disponibile in oltre 100 lingue, tra cui l'italiano) a otto emozioni di base (rabbia, paura, 
 anticipazione, fiducia, sorpresa, tristezza, gioia, disgusto) del modello psicologico di Plutchik tramite valori binari e a due polarità (positivo/negativo), 
 consentendo l'analisi delle emozioni in testi.

 \Fonte{NRC-EmoLex}


 \subsection{NRC VAD Lexicon}
Sviluppato dal National Research Council Canada, è un dizionario che fornisce punteggi continui (annotati da esseri umani) per oltre 20.000 termini lungo 
i tre assi fondamentali dell'emozione: Valenza, Arousal e Dominanza. A differenza dei modelli categoriali, permette una misurazione dimensionale diretta, 
rivelandosi strutturalmente affine allo spazio vettoriale del software.

\Fonte{NRC-VAD}


\subsection{lakh-midi}
Il dataset MIDI di Lakh è una raccolta di 176.581 file MIDI unici, di cui 45.129 sono stati abbinati e allineati alle voci del Million Song Dataset. Il suo obiettivo è facilitare il recupero su larga scala delle informazioni musicali, sia simbolico (usando solo i file MIDI) sia basato su contenuti audio (utilizzando informazioni estratte dai file MIDI come annotazioni per i file audio abbinati).

\Fonte{raffel-lmd}

Nel contesto di questo progetto il dataset lakh-midi si occupa della parte della generazione del file MIDI.
\section{Librerie}
\subsection{Pythorch}
PyTorch è una libreria tensoriale ottimizzata per il deep learning utilizzando GPU e CPU.
\Fonte{pytorch-sort}
PyTorch viene utilizzato per caricare il modello pre-addestrato, gestire la memoria e campionare le nuove note musicali in fase di inferenza.
Nello specifico è possibile vedere come torch.nn sia usato per la costruzione dell'architettura neurale.
Si farà inoltre uso di torch.softmax / F.softmax che permettono l'uso della funzione di attivazione omonima che converte le uscite grezze dal modello in una distribuzione di probabilità valida compresa tra 0 e 1.

\subsection{Pyphen}
Pyphen è una libreria Python dedicata alla sillabazione automatica e alla ricerca dei punti di divisione delle parole, basata su dizionari di 
sillabazione compatibili con il formato utilizzato da Hunspell. La libreria mette inoltre a disposizione dizionari per numerose lingue, tra cui l'italiano.


\subsection{re}
Questa libreria offre delle operazioni svolgibili sulle espressioni regolari.

Un espressione regolare specifica un set di stringhe che combacia con essa. Le funzioni di questa libreria permettono di verificare se una particolare stringa combacia con una particolare espressione regolare.
Le espressioni regolari possono essere concatenate per formare nuove espressioni regolari; se A e B sono entrambe espressioni regolari, allora AB è anch'essa un'espressione regolare. In generale, se una stringa p corrisponde ad A e un'altra stringa q corrisponde a B, la stringa pq corrisponderà ad AB. Questo vale a meno che A o B contengano operazioni a bassa priorità; condizioni ai confini tra A e B; o riferimenti a gruppi numerati. 
\Fonte{python-re}

Nel codice implemetato la libreria re serve per usare la funzione re.sub che rimuove tutti i caratteri di punteggiatura da una parola mantenendo solo le lettere e i caratteri accentati.
\subsection{UnicodeData}
La libreria da accesso all'Unicode Character Database (UCD) che definisce le proprietà dei caratteri di tutti i caratteri Unicode. 

\Fonte{python-unicodedata}

Viene importata per gestire correttamente la codifica dei caratteri speciali e accentati.
In particolare viene usato il metodo lookup che cerca un carattere dal nome. Se un carattere con tale nome viene trovato ritorna il carattere corrispondente, altrimenti viene lanciato un Keyerror.


